\section{Object-Oriented Programming}
\label{sec:qna-oops}

\subsection*{Core concepts}

\question{What are the four pillars of OOP?}
\answer{Encapsulation (bundle data with methods and hide the data),
abstraction (expose \emph{what}, hide \emph{how}), inheritance (reuse and
specialise via an ``is-a'' relationship), and polymorphism (one interface,
many run-time behaviours). See Section~\ref{sec:lld-four-pillars}.}

\question{Difference between abstraction and encapsulation?}
\answer{Abstraction hides \emph{complexity/implementation} -- you see a
simple interface. Encapsulation hides \emph{data} -- fields are private and
guarded by methods. Abstraction is a design-level idea; encapsulation is
its language-level enforcement. They usually appear together but answer
different questions: ``what do I expose?'' vs.\ ``how do I protect it?''}

\question{Difference between overloading and overriding?}
\answer{Overloading is \emph{compile-time} polymorphism: same method name,
different parameter lists, resolved by the compiler. Overriding is
\emph{run-time} polymorphism: a subclass redefines a base class's
\texttt{virtual} method with the same signature, resolved by the object's
actual type via the vtable.}

\question{What is a virtual function and how is run-time polymorphism
implemented?}
\answer{A \texttt{virtual} function can be overridden and is dispatched on
the object's dynamic type. In C++ this uses a per-class \emph{vtable} (a
table of function pointers) and a per-object \emph{vptr} pointing to it;
calling a virtual function follows the vptr to the correct override at run
time.}

\question{What is a pure virtual function / abstract class?}
\answer{A pure virtual function (\texttt{= 0}) has no body and must be
overridden. A class with at least one is \emph{abstract} -- it cannot be
instantiated and serves as an interface. It is how C++ expresses ``this is
a contract, not a concrete type.''}

\subsection*{Inheritance and relationships}

\question{Compile-time vs run-time polymorphism -- which is faster and
why?}
\answer{Compile-time (overloading, templates) is faster: the call is
resolved and can be inlined at compile time, with no indirection.
Run-time (virtual calls) costs one pointer indirection through the vtable
and blocks inlining, but buys flexibility. Prefer static dispatch when the
type is known.}

\question{Why favour composition over inheritance?}
\answer{Inheritance is rigid -- it fixes behaviour at compile time, exposes
the base class's internals to subclasses, and can violate encapsulation.
Composition (has-a) lets you change behaviour at run time by swapping a
member and keeps classes loosely coupled. Use inheritance only for a true,
stable ``is-a.''}

\question{What is the diamond problem?}
\answer{In multiple inheritance, if class D inherits from B and C which
both inherit from A, D can get two copies of A's members, causing
ambiguity. C++ solves it with \emph{virtual inheritance}, which ensures a
single shared A subobject. Many languages avoid it by disallowing multiple
class inheritance.}

\question{Difference between an abstract class and an interface?}
\answer{An abstract class can hold state and concrete methods alongside
abstract ones and models an ``is-a'' with shared implementation. An
interface (in C++, a class of only pure virtual functions) is a pure
contract with no state. A class typically extends one abstract class but
implements many interfaces.}

\subsection*{Applied}

\question{What are the SOLID principles?}
\answer{Single Responsibility (one reason to change), Open/Closed (extend
without modifying), Liskov Substitution (subtypes substitutable for their
base), Interface Segregation (small focused interfaces), Dependency
Inversion (depend on abstractions). Detailed in
Section~\ref{sec:lld-solid}.}

\question{What is dependency injection and why is it useful?}
\answer{Instead of a class constructing its own dependencies, they are
passed in (via constructor or setter). This inverts control, decouples the
class from concrete implementations, and makes it testable -- you can
inject a mock. It is the practical face of the Dependency Inversion
Principle.}

\question{Shallow copy vs deep copy?}
\answer{A shallow copy duplicates the object's fields as-is, so pointer
members still point to the \emph{same} underlying data -- two objects share
it. A deep copy also duplicates what the pointers reference, giving fully
independent objects. Deep copy matters whenever an object owns heap memory
or other resources.}

\question{What is the difference between a class and an object?}
\answer{A class is a blueprint -- it defines fields and methods. An object
is a concrete instance of that blueprint, occupying memory with its own
field values. One class, many objects.}
